\documentclass[letterpaper,10.5pt]{article}
\usepackage[empty]{fullpage}
\usepackage{titlesec}
\usepackage{enumitem}
\usepackage{hyperref}
\usepackage{xcolor}
\usepackage{tabularx}

\hypersetup{colorlinks=false,pdfborder={0 0 1},pdfborderstyle={/S/U/W 1},linkbordercolor=cyan,urlbordercolor=cyan}
\setlength{\parindent}{0pt}
\setlist[itemize]{itemsep=-3pt, topsep=0pt, leftmargin=15pt}
\titleformat{\section}{\large\bfseries}{}{0em}{}[\vspace{-1pt}\titlerule\vspace{2pt}]

\begin{document}

% --- FIXED HEADER TABLE ---
\begin{tabularx}{\textwidth}{@{}X r@{}}
{\LARGE \textbf{Yohei Ohata} \hspace{10pt} Curriculum Vitae} & \\ % Added missing '&' here
\noalign{\vspace{6pt}}
Center for Information and Neural Networks & \href{mailto:yohei.ohata.ee@nict.go.jp}{yohei.ohata.ee@nict.go.jp} \\
National Institute of Information and Communications Technology & \href{https://yohei.ohata.github.io}{yohei.ohata.github.io} \\
1-4, Yamadaoka, Suita, Osaka, Japan & \\
\end{tabularx}

\vspace{10pt}

\section*{Education}
\textbf{The University of Tokyo} \hfill April 2022 -- March 2024 \\
M.S. in Human Factors Engineering \\
Department of Human and Engineered Environmental Studies \\
Cumulative GPA: 4.0/4.0 \\
{\small \textit{Thesis}: Identifying Different Thought Processes for Each Learner Using Eye-Tracking and Hidden Markov Models} \\
{\small 1st in Class (1/48), Dean’s Award, Best Department Thesis Award} \\
\par\vspace{6pt}
\textbf{Tokyo University of Science} \hfill April 2018 -- March 2022 \\
B.E. in Electrical and Electronic Engineering \\
Department of Applied Electronics \\
Cumulative GPA: 3.5/4.0 \\
{\small \textit{Thesis}: A Method of Designing a Hilbert Transformer with Transmission Zeros at Specified Positions} \\

\par\vspace{6pt}
\section*{Employment}
\textbf{National Institute of Information and Communications Technology} \hfill January 2026 -- Present \\
Research Assistant \\
Center for Information and Neural Networks (CiNet) \\
Computational Social Neuroscience Group \\
\par\vspace{3pt}
\textbf{Hitachi, Ltd.,} \hfill April 2024 -- December 2025 \\
Research Engineer \\
R\&D Group \\
 \\
\par\vspace{3pt}
\section*{Publications}
\subsection*{Articles in Peer-Reviewed Journals}
\begin{enumerate}[label=\textbf{J\arabic*}, leftmargin=35pt]
  \item \textbf{Ohata, Y.}, Hachisuka, S., Kurita, K., \& Warisawa, S. (submitted). Identifying Different Thought Processes for Each Learner Using Eye-Tracking and Hidden Markov Models. \textit{Human Factors}.
\end{enumerate}\vspace{2pt}
\subsection*{Peer-Reviewed Conference Proceedings}
\begin{enumerate}[label=\textbf{C\arabic*}, leftmargin=35pt]
  \item \textbf{Ohata, Y.}, Hachisuka, S., Kurita, K., \& Warisawa, S. (2024). Visualizing and Revealing the Difference of Learner's Thought Process and Question Solving Method. \textit{In Proceedings of International Conference on Information and Communication Technology}.
  \item \textbf{Ohata, Y.}, Takao, K., Natori, T., \& Aikawa, N (2021). A Method of Designing a Hilbert Transformer with Transmission Zeros at Specified Positions. \textit{ICICE Technical Report}.
\end{enumerate}\vspace{2pt}
\subsection*{Conference Presentations}
\begin{enumerate}[label=\textbf{P\arabic*}, leftmargin=35pt]
  \item \textbf{Ohata, Y.}, Hachisuka, S., Kurita, K., \& Warisawa, S. (2023). Analysis of How Learners Learn Using Graphic Questions. \textit{The 42nd Annual Conference of Japanese Society for Educational Technology}.
  \item \textbf{Ohata, Y.}, Hachisuka, S., Kurita, K., \& Warisawa, S. (2022). Identifying Perceptual Factors that Influence Impression of a Teacher in a Classroom. \textit{The 41st Annual Conference of Japanese Society for Educational Technology}.
\end{enumerate}\vspace{2pt}
\subsection*{Patents}
\begin{enumerate}[label=\textbf{T\arabic*}, leftmargin=35pt]
  \item \textbf{Ohata, Y.}, Jin, Y. (2024). ``Management Issues Estimation Device''. JP-Patent.
  \item \textbf{Ohata, Y.}, Jin, Y. (2025). ``Management Issues Estimation System Using Bayesian Network and Large Language Models''. JP-Patent.
\end{enumerate}\vspace{3pt}
\section*{Research Experience}
\textbf{Research Assistant} \hfill January 2026 -- Present \\ 
National Institute of Information and Communications Technology \\ 
{\small {Laboratory: Computational Neuroscience \& Decision Making Laboratory}} \\ 
{\small \textit{Supervisor: Haruno Masahiko}} \\ 
\begin{itemize}
  \item Investigating papers dealing with alignment between LLM representations and human fMRI responses using brain decoding foundation models
  \item Developed a toy attention-only Transformer model to examine how contextual representations are generated, updated, and used for next-token prediction
\end{itemize}
\vspace{8pt}
\textbf{Research Assistant} \hfill September 2025 -- Present \\ 
The University of Tokyo \\ 
{\small {Laboratory: Amano \& Nakayama Lab}} \\ 
{\small \textit{Supervisor: Yamashita Ayumu, Amano Kaoru}} \\ 
\begin{itemize}
  \item Performed denoising on the DMCC55B Open fMRI dataset (AX-CPT, Cued Task Switching, Sternberg, Stroop) using BrainSync for synchronized cross-subject temporal alignment (MATLAB).
  \item Applied NASCAR tensor decomposition to the DMCC55B dataset to decompose data into spatial, temporal, and subject-load modes (MATLAB).
  \item Quantified cognitive control performance by indexing RT and Accuracy via the Cognitive Control Index, extracting contributory components using conditional randomization testing (R).
\end{itemize}
\vspace{8pt}
\textbf{Research Assistant} \hfill April 2022 -- March 2024 \\ 
The University of Tokyo \\ 
{\small {Laboratory: Inovative Learning Creation Lab}} \\ 
{\small \textit{Supervisor: Hachisuka Satori, Warisawa Shin'ichi}} \\ 
\begin{itemize}
  \item Designed experimental protocols for both adult and elementary school participants, including the development of a web-based UI and logging system using HTML/CSS/JS.
  \item Managed high-fidelity eye-tracking data collection and processing using Tobii Pro 3 and Tobii Pro Lab.
  \item Performed advanced statistical analysis using Linear Mixed-Effects Models (LMM) in R.
  \item Modeled the temporal dynamics of scan paths using Hidden Markov Models (HMM) implemented in Python.
  \item Conceptualized research goals, authored manuscripts.
\end{itemize}
\vspace{8pt}
\textbf{Research Assistant} \hfill May 2021 -- December 2021 \\ 
Tokyo University of Science \\ 
{\small {Laboratory: Signal Processing Lab}} \\ 
{\small \textit{Supervisor: Aikawa Naoyuki}} \\ 
\begin{itemize}
  \item Proposed a Composite Hilbert Transformer that integrates a pre-normalized notch filter to suppress noise at specific frequencies.
  \item Verified significant SNR improvements over conventional Band Pass Hilbert Transformer through extensive simulations in MATLAB.
\end{itemize}
\vspace{8pt}
\section*{Professional Experience}
\textbf{Research Engineer} \hfill April 2024 -- December 2025 \\ 
Hitachi, Ltd., Digital Innovation R\&D \\ 
\begin{itemize}
  \item Developed concepts and business use cases
  \item Defined functional requirements and methodologies
  \item Authored patent applications
  \item Applied LLMs in service development
  \item Established evaluation criteria for LLM verification
\end{itemize}
\vspace{8pt}
\section*{Teaching Experience}
\textbf{Sensing and Signal Processing} \hfill April 2023 -- June 2023 \\
Teaching Assistant, The University of Tokyo \\
\begin{itemize}
  \item Assisted in building signal processing systems using LabVIEW
\end{itemize}\vspace{5pt}
\section*{Honors and Awards}
IEEE CAS JJC Best Student Award \hfill 2021 \\ 
\section*{Languages}
\begin{itemize}[leftmargin=15pt]
  \item High Level in English
  \item Native in Japanese
\end{itemize}\vspace{5pt}
\section*{References}
\textbf{Haruno Masahiko, Ph.D.,} \\ 
Professor, National Institute of Information and Communications Technology \\ 
Email: \href{mailto:mharuno@nict.go.jp}{mharuno@nict.go.jp}\vspace{5pt} \\ 
\textbf{Yamashita Ayumu, Ph.D.,} \\ 
Assistant Professor, Advanced Telecommunications Research Institute International \\ 
Email: \href{mailto:ayumu@atr.jp}{ayumu@atr.jp}\vspace{5pt} \\ 
\textbf{Hachisuka Satori, Ph.D.,} \\ 
Associate Professor, The University of Tokyo \\ 
Email: \href{mailto:hachisuka@edu.k.u-tokyo}{hachisuka@edu.k.u-tokyo}\vspace{5pt} \\ 
\end{document}